\section{Robust Gaussian target mass}\label{sec:cylinder}

This appendix supplies the target probability bound used by the robustness-based theorems, for every $N\ge C$ and every $0<t\le1$.
Take balanced clone groups of sizes $k_a$, so $N/k_a\le2C$, and let $P$
project onto group-constant vectors. Every unused circuit slot is filled by
a Boolean constant, so the ideal wire vector $b^s(x)$ lies in the range of
$P$ and has norm $\sqrt N$.
Put
\[
 e_*=(16\sqrt{2C})^{-1},\quad \tau_0=e_*/4,\quad
 b_L=\sqrt{16\log9+8\log(4L)},
\]
\[
 A=\max\{32b_L,2\log(8/e_*)\},\qquad
 \rho=\min\{1/4,A/[4(d+1)],Ae_*/(64\sqrt C)\}.
\]
Use tanh gate gain $A$ and readout coefficient $2$. The gain need not be
bounded by logical circuit width; it can be accommodated by polynomial
padding if invoking the ``weights bounded by width'' complexity convention.

Allow any $0<t\le1$. The standardized coordinates
$\sqrt d W^1$, $\sqrt N W^s$, and $Nv$ have common standard deviation $t$;
the exponent-$\gamma$ case is $t=N^{(1-\gamma)/2}$.
Restrict every first-layer coordinate, every
internal row's projection onto each normalized clone-group direction, and
every readout coordinate to an interval of radius $\rho t$ around its clone
value. Then
\[
 W^s=W_*^s+R^s+E^s,\qquad R^s=R^sP,
 \quad\|R^s\|_\op\le\rho t\sqrt C,\quad E^sP=0.
\]
The free Gaussian bulk $E^s$ is independent of the cylinder. A $1/4$-net
Gaussian operator bound makes $\|E^s\|_\op\le b_Lt$ simultaneously with
conditional probability at least $1/2$.

We check robustness to a separate per-layer RMS budget $\tau_0$, which is
stronger than the total-budget requirement for $r\le\tau_0$. Write
$e_s=\|h^s_\eta-b^s\|_2/\sqrt N$. If $e_{s-1}\le e_*$, each clone-group
mean has error at most $1/16$, so the teacher part of every intended signed
preactivation is at least $A/2$. For every real $u$ and $\epsilon\in\{-1,1\}$,
\[
 \epsilon a\ge A/2\quad\Longrightarrow\quad
 |\tanh(a+u)-\epsilon|\le2e^{-A/2}+8|u|/A.
\]
This is a global secant bound: for $|u|<A/4$ use saturation, and otherwise
use the trivial bound $2$. It handles concentrated adversarial noises.
Since $E^sb^{s-1}=0$ and noisy states have norm at most
$(1+\tau_0)\sqrt N$,
\[
 e_s\le2e^{-A/2}+(8/A)\rho\sqrt C(1+\tau_0)
 +(8b_L/A)e_{s-1}+\tau_0\le e_*.
\]
The first layer starts the invariant because its coordinatewise
preactivation error is at most $\rho(d+1)\le A/4$.
The readout has noisy signed value at least
$2-1/8-\rho(1+\tau_0)-\tau_0>1$.

There are $Nq$ constrained independent scalar coordinates, where
$q=d+2+(L-2)C$. Their squared centers sum to at most $NB_0'$, with
\[
 B_0'=9dA^2+18(L-2)CA^2+8C.
\]
The elementary Gaussian interval bound and the independent operator event give
\[
 -\log\mu_{N,t}(\mathcal R_{r,N,f})
 \le \frac{N}{t^2}\left[B_0'+q\log\frac{e\sqrt{2\pi}}{2\rho}+\log2\right]
 =DN/t^2.
\]
Here $D$ is jointly polynomially bounded in $d,C,L$. Exponentially small
coordinate precision would in any case cost only its logarithm in this
negative log-mass estimate; it should not be replaced unnecessarily by its
reciprocal.

